\documentclass[11pt]{article}
\usepackage{amsmath,amssymb,amsthm,hyperref}
\usepackage[margin=27mm]{geometry}
\newtheorem{theorem}{Theorem}
\newtheorem{lemma}[theorem]{Lemma}
\newtheorem{proposition}[theorem]{Proposition}
\newcommand{\Q}{\mathbb Q}
\newcommand{\Z}{\mathbb Z}
\newcommand{\C}{\mathbb C}
\newcommand{\hF}{h_{\mathrm F}}
\newcommand{\OK}{\mathcal O_K}
\title{Evaluation of differentials over an unramified fibre}
\author{}\date{9 November 2026}
\begin{document}\maketitle

\noindent This is a candidate proof for independent mathematical review.
The argument below does not use bounds for Green functions or for
Faltings's delta invariant.

\section{A height bound from the discriminant of one fibre}

Let $X$ be a smooth projective geometrically connected curve over a
number field, of genus $g\geq1$, and let
$\pi:X\longrightarrow\mathbb P^1$ be a finite morphism of degree $d$,
unramified outside $0,1,\infty$. Write $t$ for the coordinate on the
target. Enlarge the number field $K$ so that $X$ has semistable
reduction and every point in $\pi^{-1}(2)$ is $K$-rational.
Let $\mathcal Y$ be the normalization of $\mathbb P^1_{\OK}$ in
$K(X)$. This is finite and flat over $\mathbb P^1_{\OK}$: finiteness
follows from excellence, and flatness from normality in dimension two
and the regularity of the target. On the affine $t$-chart write
$\mathcal Y=\operatorname{Spec}B$, where $B$ is a finite projective
$\OK[t]$-algebra of rank $d$. Put
\[
 C=B/(t-2)B\subset K^d,\qquad O=\OK^d.
\]
The algebra $C$ is finite locally free over $\OK$, with generic fibre
$K^d$, and every element is integral; hence $C\subset O$.
Let $\mathfrak D(C)$ be its discriminant ideal, defined by the
determinant of its trace pairing. Denote its absolute norm by
$\mathrm N\mathfrak D(C)$.

\begin{proposition}\label{prop:evaluation}
With the Faltings metric induced by the inner product
$\frac i2\int\omega\wedge\overline\omega$ on holomorphic
differentials, one has
\[
 \hF(\operatorname{Jac}X)
 \leq\frac{\log\mathrm N\mathfrak D(C)}{2[K:\Q]}
       +g\log\frac2{\sqrt\pi}.
\]
A different standard normalization of the Faltings metric changes
only an absolute constant times $g$.
\end{proposition}

\begin{proof}
Let $\mathcal X$ be a minimal regular semistable model of $X$, and
let $E=H^0(\mathcal X,\omega_{\mathcal X/\OK})$. This is a projective
$\OK$-module of rank $g$, with its $L^2$ metric at each complex
embedding, and its normalized arithmetic degree is the stable
Faltings height of the Jacobian.

Regular differentials on $\mathcal X$ define sections of the
dualizing sheaf of $\mathcal Y$. Indeed, take a common regular model
dominating $\mathcal X$ and $\mathcal Y$. Pullback preserves global
regular differentials under blowups; a differential on the common
model is regular at every codimension-one point of the normal
surface $\mathcal Y$, and therefore extends to its reflexive
dualizing sheaf. On the affine $t$-chart, finite duality identifies
this sheaf with
\[
 \operatorname{Hom}_{\OK[t]}(B,\OK[t])\,dt.
\]
Specialization at $t=2$ commutes with this dual, because $B$ is
finite projective. It therefore gives a map
\[
 E\longrightarrow C^\vee:=\operatorname{Hom}_{\OK}(C,\OK)
       \subset K^d,
\]
where the inclusion is by the trace pairing on $K^d$. Its coordinates
are the values $\omega/dt$ at the points of $\pi^{-1}(2)$.

Riemann--Hurwitz gives
$2g-2=-2d+\deg\operatorname{Ram}(\pi)\leq d-3$, since the
ramification contribution over each of three points is at most
$d-1$. In particular $d>2g-2$. A nonzero holomorphic differential
cannot vanish at all $d$ distinct points of this fibre. Thus the
evaluation map has rank $g$. Choose a set $I$ of $g$ coordinate
positions for which the projection remains injective, and put
$F=\operatorname{pr}_I(C^\vee)\subset K^g$.

The trace pairing makes $O=\OK^d$ self-dual. Hence
\[
 O\subset C^\vee,\qquad
 \operatorname{length}_{\OK_{\mathfrak p}}(C^\vee/O)
 =\operatorname{length}_{\OK_{\mathfrak p}}(O/C)
 =\tfrac12 v_{\mathfrak p}(\mathfrak D(C))
\]
for every finite prime $\mathfrak p$. Projection yields a surjection
$C^\vee/O\twoheadrightarrow F/\OK^g$. Equip $K^g$ with its standard
Euclidean metric at every complex embedding. It follows that
\[
 \widehat{\deg}(F)
 \leq \tfrac12\log\mathrm N\mathfrak D(C).
\]

For an embedding $\sigma:K\hookrightarrow\C$, each point of the
fibre lies on a separate analytic sheet above
$D=\{t:|t-2|<1/2\}$, since $D$ is simply connected and avoids the
branch locus. If $\omega=f(t)\,dt$ on such a sheet, the mean-value
inequality gives
\[
 |f(2)|^2\leq\frac4\pi\int_D|f(t)|^2\,dx\,dy
           \leq\frac4\pi\|\omega\|_{L^2}^2.
\]
Thus each row of the $g\times g$ evaluation matrix in an orthonormal
basis has Euclidean norm at most $2/\sqrt\pi$. By Hadamard's
inequality, the induced determinant map $\det E\to\det F$ has norm
at most $(2/\sqrt\pi)^g$ at every embedding. This is an integral
nonzero map of line bundles. Taking arithmetic degrees gives
\[
 \widehat{\deg}(E)
 \leq\widehat{\deg}(F)
       +[K:\Q]g\log(2/\sqrt\pi).
\]
The preceding inequality proves the proposition.
\end{proof}

\section{The finite discriminant estimate}

Let $S$ be a finite set of rational primes. Assume that the
normalization defining the cover is finite etale over
\[
 \operatorname{Spec}\Z[S^{-1},t,t^{-1},(t-1)^{-1}]
\]
(or over the corresponding base change to its field of definition).
This property persists on normalization after any finite base
extension, since a finite etale algebra over a normal base is normal.
Write $\operatorname{rad}(S)=\prod_{p\in S}p$ and $\omega(S)=\#S$.

\begin{lemma}\label{lem:discriminant}
One has
\[
 \frac{\log\mathrm N\mathfrak D(C)}{[K:\Q]}
 \leq d\left(\log2+\log\operatorname{rad}(S)
              +\omega(S)\log d\right).
\]
\end{lemma}

\begin{proof}
The trace determinant of $B/\OK[t]$ is a section of an invertible
sheaf, and defines its effective discriminant divisor. Base change
to $t=2$ identifies its intersection with this section with
$\log\mathrm N\mathfrak D(C)$. At a codimension-one point of the
base, its coefficient is $\sum_j f_jr_j$, where $e_j,f_j,r_j$ are
the ramification index, residue degree, and different exponent of
the extensions of discrete valuation rings in the normalization.
Here $\sum_j e_jf_j=d$.

At a horizontal branch divisor the residue characteristic is zero,
so $r_j=e_j-1$ and the coefficient is at most $d$. The section
$t=2$ meets $t=0$ with normalized intersection $\log2$, and meets
neither $t=1$ nor $t=\infty$. The horizontal contribution is thus
at most $d\log2$.

Let $\mathfrak p\mid p$ be a prime of $K$ and let
$e_{\mathfrak p}=v_{\mathfrak p}(p)$. Lenstra's different bound,
in the form of Javanpeykar, Proposition~4.1.3, gives
\[
 r_j\leq e_j-1+e_j e_{\mathfrak p}
                 \left\lfloor\frac{\log d}{\log p}\right\rfloor.
\]
Consequently the discriminant coefficient at the vertical divisor
over $\mathfrak p$ is at most
\[
 d\left(1+e_{\mathfrak p}
                 \left\lfloor\frac{\log d}{\log p}\right\rfloor\right).
\]
Only $p\in S$ can contribute. Sum the intersections of these
vertical divisors with the section, use
$\sum_{\mathfrak p\mid p}e_{\mathfrak p}f_{\mathfrak p}=[K:\Q]$
and $\sum_{\mathfrak p\mid p}f_{\mathfrak p}\leq[K:\Q]$, and then
use $\lfloor\log d/\log p\rfloor\log p\leq\log d$. This gives
the stated bound.
\end{proof}

\begin{theorem}\label{thm:coverheight}
Under the hypotheses of Lemma~\ref{lem:discriminant},
\[
 \hF(\operatorname{Jac}X)
 \leq\frac d2\left(\log2+\log\operatorname{rad}(S)
              +\omega(S)\log d\right)
       +g\log\frac2{\sqrt\pi}.
\]
\end{theorem}

\section{Integral modular covers}

Let $H\subset(\Z/N\Z)^\times$ contain $-1$, and let
$X_H=X_1(N)/H$. The map $t=j/1728$ is unramified outside
$0,1,\infty$. It has an integral finite etale model away from
these three sections and the primes dividing $6N$. An explicit
verification uses the elliptic curve
\[
 E_t:\quad y^2=x^3-\frac{3t}{t-1}x-\frac{2t}{t-1}
\]
over $\Z[1/(6N),t,t^{-1},(t-1)^{-1}]$. Its discriminant is
$1728t^2/(t-1)^3$ and its $j$-invariant is $1728t$.
The scheme of points of exact order $N$ is finite etale. Its
quotient by the constant group $H$ is also finite etale, because
$H$ acts freely. Its generic fibre is precisely the modular cover
$X_H\to\mathbb P^1_j$ restricted to this open set: the only
automorphisms of the elliptic curve are $\{1,-1\}\subset H$,
and hence taking an $H$-orbit removes the ambiguity in choosing an
isomorphism to $E_t$. The construction agrees with the
normalization, since this finite etale model is normal.

Theorem~\ref{thm:coverheight} therefore bounds the stable Jacobian
height uniformly in $H$ by
\[
 \hF(\operatorname{Jac}X_H)
 \leq\frac d2\left(\log2+\log\operatorname{rad}(6N)
              +\omega(6N)\log d\right)+g\log\frac2{\sqrt\pi},
 \qquad d=\deg(j|_{X_H}).
\]

\section{Application to the inherited five-contact realization}

The inherited realization gives an abelian surface $A$, a squarefree
integer $R=B_*M_5(a)$, a level $N$, and an even nebentypus character
$\varepsilon$ with $H=\ker\varepsilon$, such that
\[
 \begin{gathered}
 h(a)\leq2\max(0,\hF(A))+128,\qquad
 N\leq210^{300}R^2,\quad\operatorname{rad}(N)\mid R,\\
 t=\operatorname{ord}\varepsilon\leq4608\varphi(R),\qquad
 d=\deg(j|_{X_H})
   =Nt\prod_{p\mid N}(1+p^{-1})
   \leq4608\cdot210^{300}R^3.
 \end{gathered}
\]
The curve has $g\leq d$, and $A$ is an isogeny factor of its
Jacobian over a field of degree at most $2304R$. The alternative
complex multiplication case has $\hF(A)<854R$.
The exact identity
\[
 \varphi(R)\prod_{p\mid R}(1+p^{-1})
 =R\prod_{p\mid R}(1-p^{-2})<R
\]
is used in the bound for $d$. Since $6\mid R$,
Theorem~\ref{thm:coverheight} gives
\[
 \hF(\operatorname{Jac}X_H)\ll R^3(1+\log R)^2.
\]
The stable subvariety inequality of R\'emond, followed by the
fixed-dimension isogeny bound of Gaudron--R\'emond applied to an
abelian surface factor, gives
\[
 \hF(A)\ll R^3(1+\log R)^2.
\]
The latter step only adds $O(g+\log R+\log^+\hF(\operatorname{Jac}X_H))$;
the isogeny theorem is applied in dimension two, not in dimension
$g$. Thus the candidate conclusion is
\[
 h(a)\ll M_5(a)^3(1+\log M_5(a))^2,
\]
with an absolute constant depending on the fixed bad-prime product
$B_*$. In particular, for each fixed $\eta>0$ one obtains
$h(a)\ll_\eta M_5(a)^{3+\eta}$.

\section{Primary sources to verify in the independent review}
\begin{itemize}
\item A.~Javanpeykar, \emph{Polynomial bounds for Arakelov invariants
of Belyi curves}, Algebra Number Theory 8 (2014), 89--140,
especially Section~2.3 (normalization of the Jacobian height),
Proposition~4.1.3 (different bound), and Section~4.2 (finite duality
and the discriminant divisor). DOI 10.2140/ant.2014.8.89.
\item Standard finite duality identifies the relative dualizing
module of a finite flat algebra $B/A$ with
$\operatorname{Hom}_A(B,A)$. The argument uses this module and
does not require that the dualizing module of the normal surface
be invertible.
\item G.~R\'emond, \emph{Propri\'et\'es de la hauteur de Faltings},
Ann. Sc. Norm. Super. Pisa Cl. Sci. 23 (2022), 1589--1596,
Corollary~1.2 and its stable form.
\item E.~Gaudron and G.~R\'emond, \emph{Nouveaux th\'eor\`emes
d'isog\'enie}, M\'em. Soc. Math. France 176 (2023),
Theorem~1.9(1), applied in dimension two.
\end{itemize}
\end{document}
